\begingroup\renewcommand{\arraystretch}{1.3}
\begin{longtable}{|>{\raggedright\arraybackslash}p{2.1cm}|>{\centering\arraybackslash}p{1.4cm}|>{\raggedright\arraybackslash}p{10.9cm}|}
\hline
\cellcolor{navy}\textcolor{white}{\textbf{Date}} & \cellcolor{navy}\textcolor{white}{\textbf{Version}} & \cellcolor{navy}\textcolor{white}{\textbf{Description}} \\ \hline
\endfirsthead
\hline
\cellcolor{navy}\textcolor{white}{\textbf{Date}} & \cellcolor{navy}\textcolor{white}{\textbf{Version}} & \cellcolor{navy}\textcolor{white}{\textbf{Description}} \\ \hline
\endhead
2026-09-15 & 1 & Compiled: 27+ named approaches/systems across 5 categories. \\ \hline
\rowcolor{light}
2026-10-07 & 2 & Header re-aligned to the current Main Research Question. \\ \hline
2026-10-07 & 3 & Restructured to the assignment's standard background-research template (Introduction / Background / Relevant results / Inspiring methodologies / Conclusion / References) and extended with an evaluation-methodology subsection. \\ \hline
\rowcolor{light}
2026-10-07 & 4 & Extended with Section 3.7 (systematic evaluation methodology as its own field of study) and a revised Section 4.3, after feedback that the evaluation-methodology content read as an invented project-specific protocol rather than literature-grounded research. \\ \hline
2026-10-07 & 5 & Rewritten to first person so the document reads as my own work. \\ \hline
\rowcolor{light}
2026-10-07 & 6 & Extended with Section 3.8, analysing my own InnoBeweegLab stakeholder survey -- the first primary field data in this document rather than literature. \\ \hline
2026-10-07 & 7 & Extended with the full set of DOT sub-questions (Library, Field, Lab/SQ3, Lab/SQ4, Workshop, Showroom) placed directly under the Main Research Question in Section 1.2. \\ \hline
\rowcolor{light}
2026-10-09 & 8 & Converted into a standalone LaTeX project (main.tex plus one file per chapter under sections/); renamed "Document history" to "Version history" and reformatted it as this table. \\ \hline
2026-10-09 & 9 & Removed every reference to internal working-file names -- these were never part of the submission and have no meaning to an external reader. Replaced each with either the actual figures involved or a reference to a real submitted deliverable (the Project Plan, the Requirements Engineering document). \\ \hline
\rowcolor{light}
2026-10-09 & 10 & Verified every citation against its primary source and corrected five errors: re-attributed the AIS framework to its actual source (Rashkin et al., 2021) rather than a different, related Google paper; split the GATOS Workflow citation into its three genuinely distinct sources (Katz et al.'s original paper, its Nature HSSC 2026 peer-reviewed version, and a separate EDM 2025 paper applying it to real data) after confirming they had been conflated; completed the Chalmers/Volvo master's thesis citation with its real title, authors and institution now that the source document was available for direct verification; and corrected the title, year and author list of the "Judge's Verdict" citation. One previously-flagged citation (AlgorithmWatch, $\sim$30\% figure) was re-checked and confirmed accurate as originally written. \\ \hline
2026-10-12 & 11 & Added two new sources: PublicAgent (Montazeri, Feng \& Sha, 2026), a multi-agent open-data-analysis framework with ablation evidence on when agent specialisation helps, added to Section 3.4's architecture evidence; and AutoUAT/Test Flow (Ferreira, Viegas, Faria \& Lima, 2025), an industrial case study on LLM-generated acceptance tests with mandatory human review, added to Section 3.5's human-in-the-loop evidence. \\ \hline
\rowcolor{light}
2026-10-25 & 12 & Added Section~\ref{sec:process}, which documents how InnoBeweegLab builds a Wijkscan report today. It covers the method, what each of the three perspectives produces and in which files, how the report is assembled, where the time and the judgement go, and which report elements a tool can trace. It draws on the reference report and its dataset, a synthetic Buurtscan dataset, two documents of writing rules, my meeting notes and my survey. Added the literature on reading maps, charts and infographics automatically (Section~\ref{sec:res-maps}) and the method adopted for analysing the map images the user supplies (Section~\ref{sec:m-maps}). Added a fourth research-gap point and four definitions. Added fourteen reference entries, each checked against its primary source, and three primary sources. Former Sections 2--5 are now 3--6. Cross-references are now generated automatically, which also corrected the section given for PublicAgent in the reference list (the multi-agent pipelines subsection, not the generation-methods one). \\ \hline
\end{longtable}
\endgroup